\documentclass[11pt,letterpaper]{article}
\usepackage[margin=1in]{geometry}
\usepackage[T1]{fontenc}
\usepackage{lmodern,microtype}
\usepackage{amsmath,amssymb,amsthm,mathtools}
\usepackage{booktabs,longtable,array}
\usepackage[dvipsnames]{xcolor}
\usepackage{xurl}
\usepackage[colorlinks=true,linkcolor=MidnightBlue,citecolor=MidnightBlue,urlcolor=MidnightBlue]{hyperref}
\usepackage{enumitem,fancyhdr}
\setlist{itemsep=3pt,topsep=5pt}
\setlength{\emergencystretch}{2em}
\setlength{\headheight}{14pt}
\pagestyle{fancy}\fancyhf{}
\fancyhead[L]{\small Synchronized certificates for ZFC refutability}
\fancyhead[R]{\small ProveIt research}
\fancyfoot[C]{\thepage}
\newtheorem{theorem}{Theorem}[section]
\newtheorem{lemma}[theorem]{Lemma}
\newtheorem{proposition}[theorem]{Proposition}
\newtheorem{corollary}[theorem]{Corollary}
\theoremstyle{definition}\newtheorem{definition}[theorem]{Definition}
\theoremstyle{remark}\newtheorem{remark}[theorem]{Remark}
\newcommand{\N}{\mathbb N}
\newcommand{\Z}{\mathbb Z}
\newcommand{\Q}{\mathbb Q}
\newcommand{\Con}{\operatorname{Con}}
\newcommand{\ZFC}{\mathrm{ZFC}}
\newcommand{\ZF}{\mathrm{ZF}}
\newcommand{\Prf}{\operatorname{Prf}}
\newcommand{\file}[1]{\texttt{\nolinkurl{#1}}}
\newcommand{\pn}{P_{522}}
\newcommand{\pz}{P^{\Z}_{2088}}
\newcommand{\norm}[3]{\mathcal N(#1,#2,#3)}
\DeclareMathOperator{\lcm}{lcm}
\title{\vspace{-1em}\textbf{Synchronized arithmetic certificates\\for ZFC refutability}\\[0.5em]
\large An explicit quartic in 522 nonnegative integers,\\and a quartic in 2,088 unrestricted integers}
\author{Prepared for Vladimir Reshetnikov's ProveIt project\\\small AI-assisted mathematical research and reproducible algebra}
\date{24 October 2026}
\begin{document}
\maketitle
\begin{abstract}
We give a fully materialized polynomial $\pn\in\Z[x_0,\ldots,x_{521}]$
of exact degree four such that
\[
 \ZFC\vdash
 \bigl(\nexists x\in\N^{522}:\pn(x)=0\bigr)
 \longleftrightarrow \Con(\ZFC),\qquad \N=\{0,1,2,\ldots\}.
\]
It is the sum of squares of 400 explicitly listed quadratic residuals,
with 20,225 nonzero expanded monomials and coefficient height
83,900,628. There are no unbound coefficient ports, variable exponents,
factorial macros, or bounded quantifiers in the delivered polynomial.
An accompanying polynomial $\pz$ has 2,088 unrestricted integer
variables and exact degree four, with the same consistency equivalence.

The construction directly encodes the calibrated numerical calculus
$C_{222}$; it does not simulate K215 or invoke Jones's universal equation.
Its main new step is to synchronize arithmetic certificates across an
entire finite circuit: one Pell frame handles every exponentiation, one
auxiliary Pell solution handles every index, and one delayed radix handles
every binomial coefficient. A circuit majorant replaces a large duplicate
evaluation, and certified polynomial substitutions remove unnecessary
coordinates. Against the newest inspected direct-certificate archive,
the natural-variable count falls from 3,614 to 522, a reduction of
85.56\%, while the maximum coefficient also decreases.

These are ordinary mathematical proofs formalizable in ZFC, accompanied
by exact finite algebra audits. We do not claim global minimality,
worldwide priority, a proof of consistency, or a newly compiled
end-to-end Lean or Rocq certificate.
\end{abstract}

\section{Result, comparison, and exact deliverable}

The comparison snapshot is ProveIt commit
\file{a88aefd7ddab4771534bb049c67350e397c91b9a}. Its active research
notes and incoming archives contain stronger results than the numbered
collection's opening summary. In particular, the incoming
\file{ConZFC_Shared_Certificates_2026-10-24.zip} already supplies a direct
quartic in 3,614 natural variables. That archive, rather than the older
16,036-variable direct construction, is our immediate baseline
\cite{repository,latest,shared}.

\begin{theorem}[Delivered equations]\label{thm:main}
Let $\pn$ be the literal expanded polynomial in
\file{equations/natural522/expanded_polynomial.json}, with the coordinate
ordering in the accompanying \file{equation.json}. Then
\begin{equation}\label{eq:main}
 \ZFC\vdash
 \left(\exists x\in\N^{522}\;\pn(x)=0\right)
 \longleftrightarrow
 \left(\exists p\;\Prf_{\ZFC}(p,\bot)\right).
\end{equation}
The polynomial has exact degree four, 20,225 monomials, and coefficient
height 83,900,628. Its compact, literal definition is
\begin{equation}\label{eq:sos}
 \boxed{\displaystyle \pn(x)=\sum_{j=0}^{399}R_j(x)^2=0,\qquad x\in\N^{522},}
\end{equation}
where every $R_j$ is explicitly supplied and has degree at most two.
The 529- and 596-coordinate companion quartics have the same equivalence.
The polynomial $\pz$ in \file{equations/integer2088} satisfies
\eqref{eq:main} with $x\in\Z^{2088}$ in place of $x\in\N^{522}$.
Taking negations gives the requested no-solution/consistency equivalence.
\end{theorem}

The proof occupies the following sections. The numerical measurements are
obtained by exact expansion and coefficient collection; they are not
estimates from the syntax of an arithmetic circuit.

\begin{table}[ht]
\centering\small
\begin{tabular}{lrrrr}
\toprule
Natural quartic & Variables & Residuals & Monomials & Coefficient height\\
\midrule
Incoming direct baseline & 3,614 & 2,954 & 15,996 & 127,778,545\\
Synchronized source & 895 & 773 & 4,863 & 63,908,265\\
Sparse contraction & 596 & 474 & 7,263 & 63,909,160\\
Compact contraction & 529 & 407 & 10,782 & 211,184,624\\
Fewest variables delivered & \textbf{522} & \textbf{400} & 20,225 & 83,900,628\\
\bottomrule
\end{tabular}
\caption{All rows have exact degree four and quantify nonnegative integers.
The different support limits are deliberate expression-size tradeoffs.}
\label{tab:frontier}
\end{table}

The 522-coordinate form minimizes variables among the variants materialized
here. The 529-coordinate form uses only seven more coordinates and about
half as many monomials. The 596-coordinate form has a smaller coefficient
height and fewer monomials still. These are different useful choices;
no optimization theorem for the contraction order is asserted.

There are already much lower-variable equations in the repository. Its
active comparison records nine-positive and ten-integer equations, with
much larger degree, and a separate incoming norm construction has thirteen
integer variables. Older machine-derived quartics also have far fewer
variables, at the cost of large fixed coefficient payloads. Our result
improves the \emph{direct, fully expanded, modest-coefficient quartic
route}. It is not a replacement for every point of that broader
variable/degree/coefficient frontier. Jones's and Sun's reductions remain
relevant background, but neither is used in this construction
\cite{latest,jones,sun}.

\subsection{What makes the polynomial explicit}

Each residual file contains the complete ordered variable list and a
finite list of integer coefficients and monomials. A term is encoded as
\file{[coefficient,[variable_indices]]}; an empty index list denotes the
constant monomial and repeated indices denote powers. Indices start at
zero. The expanded file is the same polynomial after exact squaring,
summation, and collection. All coefficients are ordinary signed decimal
integers; for $\pn$, their magnitudes are less than $2^{27}$.

The archive includes a deterministic generator, independently implemented
checks, the full elimination maps, immediate source dependencies, and
cryptographic file ledgers. The hashes identify the supplied bytes; the
literal data and mathematical construction, not a hash alone, define the
equation. The construction never runs a consistency-deciding search and
never uses an answer to the consistency question in choosing a coefficient.

\section{Logical source: a numerical calculus rather than a machine}
\label{sec:calculus}

The inherited source is the numerical calculus $C_{222}$. Its false target
is code $1$, and its freely available truth label is code $0$. The
repository supplies unconditional calibration statements
\begin{equation}\label{eq:calibration}
 C_{222}(1) \longleftrightarrow\
 \operatorname{BProv}(\ZF,\bot)
 \ \longleftrightarrow\
 \operatorname{BProv}(\ZFC,\bot).
\end{equation}
The named sources are \file{KEndpoint.c222_iff_zf_refutation} and
\file{ConstructibleChoice.zf_refutation_iff_zfc_refutation}; the latter
uses the constructible-universe argument for Choice. Their immediate
source files and pinned hashes are included \cite{calibration,choice}.
The first theorem is independent of the machine parameters that occur
later in its source module. Some imported modules also support machine
theorems; we make no claim that the entire import graph is machine-free.

For completeness, here is the entire arithmetic interface being encoded.
Write
\begin{align*}
 I(a,b)&=2(a+b)(a+b+1)+4b,\\
 \mathsf M(a,b)&=I(a,b)+1,&
 \mathsf A(a,b)&=I(a,b)+2,&
 \mathsf E(a,b)&=I(a,b)+3,\\
 \mathsf N(a)&=I(a,1),&
 \mathsf X(a,b)&=\mathsf N(\mathsf A(a,\mathsf N(b))).
\end{align*}
The eight unconditional axiom codes are
\begin{align*}
 c_1&=\mathsf X(a,\mathsf E(a,b)),\\
 c_2&=I(\mathsf A(a,I(b,c)),I(\mathsf A(a,b),\mathsf A(a,c))),\\
 c_3&=I(\mathsf E(a,b),I(\mathsf M(c,a),\mathsf M(c,b))),\\
 c_4&=I(\mathsf A(a,\mathsf A(b,c)),\mathsf A(b,\mathsf A(a,c))),\\
 c_5&=I(\mathsf E(c,a),I(\mathsf M(c,b),\mathsf M(a,b))),\\
 c_6&=I(\mathsf X(a,\mathsf A(a,b)),b),\\
 c_7&=I(\mathsf E(c,a),I(\mathsf E(c,b),\mathsf A(\mathsf M(a,b),\mathsf E(a,b)))),\\
 c_8&=I(I(I(c,d),a),I(I(a,c),I(b,c))).
\end{align*}
The ninth, guarded bundle code is a fixed composition. To avoid an
unreadable nested expression, define the following polynomial abbreviations:
\begin{align*}
 \mathsf R(b,k,h)&=\mathsf A(b,I(h,\mathsf N(\mathsf A(b,I(h,\mathsf N(\mathsf M(b,k))))))),\\
 \mathsf S(a,b,k)&=\mathsf A(k,I(\mathsf M(k,a),\mathsf M(k,b))),\\
 C_0(a,b,k,h)&=I(\mathsf S(a,b,k),\mathsf R(b,k,h)),\\
 C_1(a,b,k,h)&=\mathsf N(\mathsf A(a,C_0(a,b,k,h))),\\
 C_2(a,b,k,h)&=\mathsf N(\mathsf A(a,I(h,C_1(a,b,k,h)))),\\
 \mathsf C(a,b,k,h)&=\mathsf A(b,I(\mathsf M(b,k),C_2(a,b,k,h))),\\
 \mathsf V(a,b,c)&=I(\mathsf M(a,b),\mathsf A(c,\mathsf M(a,b))),\\
 \mathsf K(a,b,c)&=\mathsf A(c,I(\mathsf C(a,b,c,c),\mathsf N(\mathsf R(b,c,c)))),\\
 c_*&=\mathsf N(I(\mathsf V(a,b,c),\mathsf K(a,b,c))).
\end{align*}
Its side condition is represented by the single equation
\[
 g=(c-a)^2(c-b)^2-(d+1)^2=0.
\]
Over $\N$, existence of $d$ is equivalent to $c\ne a$ and $c\ne b$:
take $d=|(c-a)(c-b)|-1$ in the converse direction.

There are twelve padded two-premise branches: free truth; the eight
unconditional axioms; the guarded bundle; modus ponens $q=I(p,v)$;
and generalization $q=0$, $v=\mathsf A(a,p)$. In every axiom branch
$p=q=0$. The four schema coordinates $a,b,c,d$ need satisfy only the
selected branch. Exact compositional checks match all 56 inherited wire
definitions and all 34 inherited branch residuals against these formulas.

\subsection{Factoring the common axiom side condition}

Put
\[
 G=p^2+q^2,\qquad
 L=v\prod_{j=1}^{8}(v-c_j),\qquad
 H_*=(v-c_*)^2+g^2.
\]
The union of the ten axiom/truth branches is represented by
\begin{equation}\label{eq:factored-axiom}
 G+L^2H_*=0.
\end{equation}
The quantities $G$, $L^2$, and $H_*$ are nonnegative, so this equation forces
$p=q=0$ and either $L=0$ or $H_*=0$. These are exactly the required
unconditional or guarded alternatives. Multiplying by the unchanged
modus-ponens and generalization factors gives the full Horn-rule
polynomial
\begin{equation}\label{eq:horn}
 \mathcal H=(G+L^2H_*)\,(q-I(p,v))^2
 \bigl(q^2+(v-\mathsf A(a,p))^2\bigr).
\end{equation}
This changes polynomial values away from zero. The claimed preservation
is an exact equivalence of zero sets, not a polynomial identity.

\section{Finite proofs as bounded arithmetic paths}
\label{sec:paths}

We recall the inherited commutative-flow argument because its logical
role matters. For a two-premise rule $(p,q)\to v$, assign the vector
$e_v-e_p-e_q$ in the free abelian group on nonzero labels; put $e_0=0$.
A finite proof tree of target $1$ yields a multiset of rule instances
whose total flow is $e_1$.

Conversely, let $D$ be the least set containing truth and closed under the
rules. If $1\notin D$, sum the coordinates outside $D$. A rule whose
conclusion is outside $D$ has at least one premise outside $D$, so its
contribution is nonpositive. A rule with conclusion in $D$ also has
nonpositive contribution. A nonnegative sum of rule vectors therefore
cannot equal $e_1$, whose outside-$D$ sum is one. Only the finitely many
labels in the certificate are involved. Thus the flow certificate is
equivalent to derivability, even though its rule multiset is not ordered.

\subsection{A polynomial weight carrier}

Let $t\ge1$ be the number of rule occurrences, and choose $m\ge1$ to
bound all seven local Horn coordinates $p,q,v,a,b,c,d$. Set
\[
 M=(m+1)^t,\qquad w(a)=1+aM.
\]
For $t$ rule instances define
\[
 A_0(X)=\prod_{i<t}(1+v_iX),\qquad
 C_0(X)=(1+X)\prod_{i<t}(1+p_iX)(1+q_iX).
\]
If $A_0(M)=C_0(M)$, then $\deg C_0\le t$: otherwise
$C_0(M)\ge M^{t+1}$ whereas
$(1+mM)^t<((m+1)M)^t=M^{t+1}$. Both polynomials now have coefficient
sum at most $(m+1)^t=M$, and constant coefficient one, so every
coefficient is strictly smaller than $M$. Uniqueness of base-$M$
expansion gives $A_0=C_0$. Unique factorization in $\Q[X]$ gives precisely
the required flow vector. Numerical coprimality of the weights is not needed.

It suffices to encode an accumulator path
\begin{equation}\label{eq:path}
 S_0=M+1,\quad S_t=1,\qquad
 S_{i+1}(1+v_iM)=S_i(1+p_iM)(1+q_iM).
\end{equation}
All states are positive by induction. Multiplication and cancellation
yield $A_0(M)=C_0(M)$. Conversely, expand a finite derivation into a tree
and traverse it by replacing a frontier conclusion by its premises.
The product of frontier weights supplies the required natural states.

\subsection{Eleven coordinates per local row}

Encode the path by the usual beta digits
\[
 \beta(x,y,i)=x\bmod D_i,\qquad D_i=1+(i+1)y.
\]
For a row $i<t$, the eleven local coordinates and their bounds are
\begin{center}
\begin{tabular}{lll}
\toprule
Role & Coordinates & Bound\\
\midrule
Horn rule and schemas & $p,q,v,a,b,c,d$ & $m$\\
Consecutive states & $S,S'$ & $y$\\
Beta quotients & $u,u'$ & $x$\\
\bottomrule
\end{tabular}
\end{center}
The three local equalities, besides $\mathcal H=0$, are
\begin{align}
 S'(1+vM)&=S(1+pM)(1+qM),\label{eq:local-acc}\\
 x&=S+D_i u, &x&=S'+(D_i+y)u'.\label{eq:local-beta}
\end{align}
Because $S,S'\le y<D_i,D_i+y$, these are exact remainder equations.
The two former remainder-slack coordinates are unnecessary. The quotient
bounds are valid since $x=S+D_i u$ with $D_i\ge1$ implies $u\le x$,
and similarly for $u'$.

Write $M=M_0+1$. The outer endpoint equations are
\begin{align}
 x&=M_0+2+(y+1)u_0,& y&=M_0+2+r_0,\label{eq:endfirst}\\
 x&=1+(1+(t+1)y)u_t.\label{eq:endlast}
\end{align}
Thus $y\ge M+1\ge2$, so the initial state is below its modulus and the
final state $1$ is below its modulus. An additional final remainder
slack is unnecessary.

Completeness has an essential order of choices. First take the finite
Horn proof and choose $m$ to bound all seven proof coordinates. Then
choose $M$ and construct its finite state sequence. Choose $y$ above
all states and divisible by $t!$. The moduli $D_0,\ldots,D_t$ are
pairwise coprime: a common divisor for $D_i,D_j$ divides $j-i$ and is
coprime to $y$, whereas $j-i\mid t!\mid y$. CRT supplies a natural $x$
with every required state as its digit; taking a sufficiently large
representative gives natural quotients. There is no independent uniform
local bound and no circular requirement that $y$ exceed its own code $x$.

Let $P_{\mathrm{loc}}$ be \eqref{eq:horn} plus the squares of the three
residuals in \eqref{eq:local-acc}--\eqref{eq:local-beta}. Then it is
nonnegative, and zero precisely for a valid local row. Its large degree
is an intermediate property; we subsequently represent its finite
arithmetic circuit by quadratic gate equations.

\section{Bounded universal checking and a cheap majorant}
\label{sec:crt}

A signed polynomial circuit can be split as $P=E-F$, with $E,F$
having nonnegative coefficients. Recursively use
\[
 (E_1,F_1)(E_2,F_2)
 =(E_1E_2+F_1F_2,\ E_1F_2+F_1E_2),
\]
with componentwise addition and swapped components for negation.
Identical subexpressions are shared. Apply this to $P_{\mathrm{loc}}$.
Let $z=(t,x,y,m,M_0)$, and let $\mathbf u$ be the eleven bounds above.

\begin{lemma}[Shared bounded-row certificate]\label{lem:crt}
Suppose $H$ is a natural integer satisfying
\[
 H\ge E(t,z,\mathbf u)+F(t,z,\mathbf u)+t+x+y+m+1.
\]
Put $T=H!$ and $\mathcal B=\prod_{i<t}(1+(i+1)T)$. Then
\[
 \forall i<t\ \exists w\le\mathbf u:\ E(i,z,w)=F(i,z,w)
\]
is equivalent to existence of natural $C,A_1,\ldots,A_{11}$ satisfying
\begin{align}
 \mathcal B&\mid\binom{A_j}{u_j+1} &&(1\le j\le11),\label{eq:range}\\
 \mathcal B&\mid 1+(C+1)T,\label{eq:crtindex}\\
 E(C,z,A)&\equiv F(C,z,A)\pmod{\mathcal B}.\label{eq:crteval}
\end{align}
All $A_j$ may be required strictly positive.
\end{lemma}
\begin{proof}
Let $d_i=1+(i+1)T$. The $d_i$ are pairwise coprime by the same gcd
argument as for beta digits, and each is coprime to $T$. Every prime
$\ell\mid d_i$ exceeds $H$, since a prime at most $H$ divides $H!$.
The identity
\[
 (u_j+1)!\binom{A_j}{u_j+1}=\prod_{r=0}^{u_j}(A_j-r)
\]
shows from \eqref{eq:range} that $A_j$ is congruent modulo $\ell$ to
some $w_j\in\{0,\ldots,u_j\}$. From \eqref{eq:crtindex}, cancellation
of $T$ gives $C\equiv i\pmod\ell$. Equation \eqref{eq:crteval} then
gives equality of $E(i,z,w)$ and $F(i,z,w)$ modulo $\ell$.
Both are nonnegative and smaller than $\ell$, by monotonicity and the
bound on $H$, hence they are equal as integers.

Conversely, choose a witness for each of the finitely many rows and
apply CRT modulo the pairwise coprime $d_i$ to encode its coordinates
and the index. A vanishing falling-factorial factor gives
$\mathcal B\mid (u_j+1)!\binom{A_j}{u_j+1}$. Since
$(u_j+1)!\mid H!$ and $\gcd(\mathcal B,H!)=1$, cancel the factorial.
The other two congruences follow directly. Adding positive multiples of
$\mathcal B$ makes all coordinate representatives positive without
affecting any of these arguments.
\end{proof}

\subsection{Replacing a duplicate circuit by 23 squares}

The baseline evaluates a large copy of $E+F$ at the bounding inputs to
construct $H$. We need only an upper bound. Put
\[
 R_0=2+t+x+y+m+M_0.
\]
Every input is at most $R_0$, and $R_0\ge2$. Give each input exponent
one and each constant $c\ge0$ an exponent $e(c)$ with $c\le2^{e(c)}$.
Use $e(0)=e(1)=0$, and otherwise $e(c)=\lceil\log_2c\rceil$.
At a sum gate use $\max(e_1,e_2)+1$; at a product gate use $e_1+e_2$.
Induction proves that every node is at most $R_0$ to its assigned
exponent. This induction applies to the actual nonnegative pair circuit,
not to a signed polynomial after cancellation.

The supplied majorant certificate checks 388 gate inequalities for
$E+F+t+x+y+m+1$ and gives exponent $6,893,159$. Therefore the convenient
choice
\begin{equation}\label{eq:majorant}
 \boxed{H=R_0^{8,388,608}=R_0^{2^{23}}}
\end{equation}
is sufficient and uses only 23 squarings. This is a fixed polynomial
power, subsequently represented by 23 quadratic defining equations.
It is not a variable exponent left inside the final equation.

\section{One radix for every binomial coefficient}
\label{sec:radix}

The bounded certificate still contains binomial coefficients. The direct
compiler has fourteen such requests, including those used to represent
the two factorials and the modulus product. Treat their inputs and outputs
as a finite semantic arithmetic circuit.

\begin{theorem}[Simultaneous binomial digit certificate]\label{thm:radix}
For a finite family of positive upper arguments $n_j$ and positive lower
arguments $k_j$, set
\[
 N=\sum_{n\in\{n_j\}\text{ as distinct circuit atoms}} n,
 \qquad q=2^N+1.
\]
For each request introduce $C_j,L_j,U_j,h_j,r_j,s_j,d_j\in\N$ and impose
\begin{align}
 L_j&=q^{k_j},&U_j&=(q+1)^{n_j},\label{eq:digitpowers}\\
 U_j&=L_j(qh_j+C_j)+r_j,\label{eq:digitmain}\\
 L_j&=r_j+s_j+1,&q&=C_j+d_j+1.\label{eq:digitbounds}
\end{align}
Then these equations are equivalent to $C_j=\binom{n_j}{k_j}$ for
every $j$. Equal power requests are represented once.
\end{theorem}
\begin{proof}
Each $n_j\le N$, so every binomial coefficient of $(1+X)^{n_j}$ is
at most $2^{n_j}<q$. Hence $(q+1)^{n_j}$ has a carry-free base-$q$
expansion with exactly those coefficients as digits. Equations
\eqref{eq:digitmain}--\eqref{eq:digitbounds} isolate digit $k_j$.
This also handles $k_j>n_j$, when the digit is zero. Conversely take
the ordinary quotient, remainder, and selected digit of that expansion;
all slack variables are natural.
\end{proof}

The crucial implementation choice is to \emph{defer} the digit certificates.
A request for $\binom nk$ first allocates its output coordinate. After
all requests are known, the generator constructs one common $q$ and the
digit equations. An upper argument can depend on another binomial's
semantic output; this does not invalidate the theorem, which holds
simultaneously for arbitrary natural assignments to all its arguments.
The auxiliary radix is allowed to be chosen after the semantic values.

In the emitted system, fourteen binomial requests require one common radix,
shared denominators, and shared power-index certificates. The total number
of distinct power requests in the entire compiler drops from 55 in the
baseline to 27. The next section certifies these 27 requests together.

\section{One Pell frame for an entire power circuit}
\label{sec:power}

For $a\ge2$, define the Pell sequences by
\begin{align*}
 X_0(a)&=1,&Y_0(a)&=0,\\
 X_{j+1}(a)&=aX_j(a)+(a^2-1)Y_j(a),&
 Y_{j+1}(a)&=X_j(a)+aY_j(a).
\end{align*}
They satisfy $X_j(a)^2-(a^2-1)Y_j(a)^2=1$. We use the standard
classification of natural solutions of this norm equation and the finite
Matiyasevich Pell-index criterion. These statements are available in the
cited Mathlib development and its repository interface
\cite{pell,carneiro,pellinterface}. The full inherited exponentiation
criterion is not needed: the shared frame supplies its own exponential bound.

\subsection{The common exponential bound}

For every distinct circuit atom occurring as a base $n$, exponent $k$, or
proposed output $m$ of a power request, define
\begin{equation}\label{eq:frame}
 W=2+\sum\text{these atoms},\qquad Z=z+1,\qquad
 A^2-W^3(W+2)Z^2=1.
\end{equation}
All coordinates in this section are natural. For each distinct base $n$
introduce a natural modulus $T_n$ with
\begin{equation}\label{eq:basemod}
 2An=T_n+n^2+1.
\end{equation}
Equation \eqref{eq:basemod} excludes $n=0$ in any solution. The
positivity of every exponent $k$ is a separate, checked syntactic property
of the source. Thus $1\le n,k\le W$ and $0\le m\le W$.

\begin{lemma}[The frame is exponentially large]\label{lem:frame}
Every solution of \eqref{eq:frame} and \eqref{eq:basemod} has
\[
 A\ge(W+1)^W>W^W,
 \qquad T_n\ge2A-2>W^W\ge n^k.
\]
In particular, $m<T_n$ for every power request. Conversely, after any
finite collection of intended power values is fixed, a common frame exists.
\end{lemma}
\begin{proof}
The norm in \eqref{eq:frame} is
$A^2-((W+1)^2-1)(WZ)^2=1$. Pell classification gives
$A=X_j(W+1)$ and $WZ=Y_j(W+1)$. Reducing the recurrences modulo $W$
gives $X_j(W+1)\equiv1$ and $Y_j(W+1)\equiv j$. Hence $W\mid j$.
Since $WZ>0$, we have $j\ge W$. The recurrence gives
$X_j(W+1)\ge(W+1)^j$, proving the exponential bound.

Also $A>W^2$ follows directly from the norm. The expression
$2An-n^2-1$ is increasing for $1\le n\le W$, and its value at one is
$2A-2$. Since $W\ge2$, this is larger than $W^W$, as required.
For completeness choose
\[
 A=X_W(W+1),\qquad Z=Y_W(W+1)/W.
\]
The congruence above makes $Z$ a positive integer and supplies the norm.
\end{proof}

The estimate $A>W^2$ by itself would not justify power recovery. The
divisibility of the Pell index by $W$, and hence the exponential estimate
in Lemma~\ref{lem:frame}, is essential.

\subsection{One auxiliary Pell solution for every index}

There are 22 distinct exponent atoms among the 27 requests. For each
one, denoted $k_i$, introduce $X_i,Y_i,S_i,R_i$ and natural slacks
$e_i,s_i,r_i$. Use one common auxiliary triple $U,V,B$ and impose
\begin{align}
 Y_i&=k_i+e_i,&X_i^2-(A^2-1)Y_i^2&=1,\label{eq:index1}\\
 S_i^2-(B^2-1)R_i^2&=1,&
 S_i&=X_i+Us_i,&R_i&=k_i+4Y_ir_i.\label{eq:index2}
\end{align}
The common constraints are
\begin{align}
 U^2-(A^2-1)V^2&=1,\label{eq:uvnorm}\\
 V&=\left(\prod_{i=1}^{22}Y_i^2\right)(v_0+1),\label{eq:uvproduct}\\
 B&=1+4V(b_0+1),&B&=A+Ub_1.\label{eq:bcommon}
\end{align}
The product in \eqref{eq:uvproduct} has the fixed length 22; it becomes
a fixed list of quadratic multiplication gates in the final polynomial.

\begin{lemma}[Simultaneous index correctness]\label{lem:index}
Constraints \eqref{eq:index1}--\eqref{eq:bcommon} force
$X_i=X_{k_i}(A)$ and $Y_i=Y_{k_i}(A)$ for every $i$. Conversely, for
any finite family of positive $k_i$ and $A\ge2$, these constraints have
simultaneous natural witnesses with those Pell coordinates.
\end{lemma}
\begin{proof}
For soundness, $Y_i\ge k_i\ge1$, $V>0$, and $Y_i^2\mid V$.
Moreover $B>1$, $B\equiv1\pmod{4Y_i}$, and $B\equiv A\pmod U$.
Together with \eqref{eq:index1}--\eqref{eq:index2}, these are exactly
the nonzero branch of the finite Matiyasevich index criterion.

For completeness, first set $X_i=X_{k_i}(A)$ and $Y_i=Y_{k_i}(A)$.
Let $L=4\prod_iY_i^2$. In the finite ring
\[
 (\Z/L\Z)[s]/(s^2-(A^2-1)),
\]
the element $A+s$ is a unit, with inverse $A-s$. A positive power is
therefore one. For its positive exponent $r$, put
$U=X_r(A)$, $V=Y_r(A)$. Then $U\equiv1\pmod L$, $V\equiv0\pmod L$,
and $V>0$. In particular $U$ is odd, and the Pell norm implies
$\gcd(U,V)=1$. Thus $\gcd(U,4V)=1$.

CRT gives $B\equiv A\pmod U$ and $B\equiv1\pmod{4V}$. Take a
sufficiently large representative $B>A$. This supplies the nonnegative
quotients in \eqref{eq:bcommon} and \eqref{eq:uvproduct}.
Set $S_i=X_{k_i}(B)$ and $R_i=Y_{k_i}(B)$. The Pell recurrences are
polynomial in the base, so $S_i\equiv X_i\pmod U$. Also
$Y_k(1)=k$, so $R_i\equiv k_i\pmod{4Y_i}$. Monotonicity in the base
gives $S_i\ge X_i$, and $Y_k(B)\ge k$ gives $R_i\ge k_i$.
All displayed quotients and slacks are therefore natural.
\end{proof}

\subsection{Recovering each power with a nonnegative quotient}

For a request with exponent $k_i$, base $n$, and output $m$, impose
one additional equation, with $h\in\N$:
\begin{equation}\label{eq:power-final}
 X_i+nY_i=AY_i+m+T_nh.
\end{equation}
For the true Pell sequences put $H_j=X_j(A)-(A-n)Y_j(A)$. Direct
substitution gives
\[
 H_0=1,\qquad H_{j+1}=nH_j+T_nY_j(A).
\]
Consequently the exact identity
\begin{equation}\label{eq:quotient-identity}
 H_k-n^k=T_n\sum_{j=0}^{k-1}n^{k-1-j}Y_j(A)
\end{equation}
has a nonnegative quotient. Soundness of \eqref{eq:power-final} now
follows from Lemmas~\ref{lem:frame} and \ref{lem:index}: $m$ and $n^k$
are congruent modulo $T_n$ and both lie in $[0,T_n)$, so they are equal.
Completeness uses the sum in \eqref{eq:quotient-identity} as $h$.

The compiler shares the index block when the exponent atom is identical,
the modulus when the base atom is identical, and the complete request when
both are identical. It does not identify distinct expressions merely
because they happen to have the same value in a particular assignment.

\subsection{Why global sharing is not circular}

For soundness, \eqref{eq:basemod} first forces every base positive.
All exponent arguments are positive by arithmetic syntax: $t=t_0+1$,
$H=R_0^{2^{23}}$, shifted factorial arguments, positive encoded CRT
coordinates, lower bounds plus one, and the sum of positive upper
arguments. The entire power network is therefore correct as a relation
on arbitrary proposed values. Next Theorem~\ref{thm:radix} identifies
all binomial outputs simultaneously. Only then is induction on the
original semantic macro circuit used to identify factorials and the
modulus product. No factorial positivity is assumed to establish the
power network: even a factorial output proposed as a zero base is
excluded by \eqref{eq:basemod} itself.

For completeness, evaluate the original finite semantic circuit on its
true inputs. Choose the common radix from those upper arguments and
evaluate its associated powers. All 27 power requests now have true,
finite values. Choose $W$ and its Pell frame; then choose the common
index certificate. Auxiliary scale choices occur after the semantic
values they bound. Their size is unrestricted, as is appropriate for an
existential Diophantine certificate.

\section{Eliminating factorials, products, and arithmetic gates}
\label{sec:materialization}

We retain the elementary arithmetic graph identities of the incoming
direct construction, with their power and binomial requests now handled
by the shared machinery.

For $n\ge1$, put
\[
 \rho=(n+2)^{n+2},\qquad c=\binom{\rho+n}{n},\qquad P=(\rho+1)^n.
\]
Then $n!=\lceil P/c\rceil$. For $n\ge2$,
\[
 \frac Pc=n!\prod_{i=1}^n\frac{\rho+1}{\rho+i},\qquad
 0\le n!-\frac Pc<\frac{n!n(n-1)}{2(\rho+1)}<1;
\]
the case $n=1$ is immediate. Hence a proposed $f=n!$ is characterized
by two natural remainder slacks and
\[
 P+c=cf+r+1,\qquad c=r+s+1.
\]
These encode $c(f-1)<P\le cf$.

The modulus product in Lemma~\ref{lem:crt} is also eliminated. Let
\[
 Q=(1+(t+1)T)^t+2,\qquad L=TQ-1.
\]
The exact identity is
\begin{equation}\label{eq:modproduct}
 \mathcal B=\left(T^t\,t!\binom{Q+t}{t}\right)\bmod L.
\end{equation}
The numerator is $\prod_{i=1}^t(TQ+iT)$ and reduces modulo $L$ to
$\prod_{i=1}^t(1+iT)$. Moreover
\[
 \mathcal B\le(1+tT)^t<(1+(t+1)T)^t+1\le L.
\]
Thus the indicated residue is the actual product. An ordinary quotient
and remainder, with a strict remainder bound, give its finite polynomial
graph. The compiler needs exactly two factorial requests and fourteen
binomial requests; all are eliminated by the preceding constructions.

Every addition or multiplication in the resulting fixed circuit can be
represented by a natural coordinate $z$ and the quadratic residual
$z-a-b$ or $z-ab$. Constants are literal numerals. A finite conjunction
of residual equations is replaced by their sum of squares. The inherited
exporter first removes 777 nonnegative affine gate coordinates, yielding
the 895-coordinate source in Table~\ref{tab:frontier}. This exporter is
included, and its output is independently regenerated and expanded.

\subsection{Certified degree-preserving contraction}

\begin{lemma}[Nonnegative polynomial substitution]\label{lem:contract}
Suppose a current residual is $\varepsilon(x-F(y))$, where
$\varepsilon\in\{-1,1\}$ and $F$ has nonnegative integer coefficients
and does not contain $x$. Delete this residual and substitute $F$ for
$x$ in every other residual. The nonnegative zero sets are in bijection,
with inverse $x=F(y)$. If $F$ is affine, quadratic residual degree is
preserved. The same holds when $F$ is quadratic and $x$ occurs only
as a degree-one monomial in every other residual.
\end{lemma}
\begin{proof}
Any original zero has $x=F(y)$, so projects to a new zero. Conversely,
evaluate $F$ on a new natural zero. Its value is natural and restores
every original equation. The degree assertions follow by direct
substitution. In the quadratic case, the pivot occurs in no product
or square elsewhere, so replacement cannot create degree above two.
\end{proof}

This is applied to literal exported residuals, including witness
equalities, not only to expressions marked as gates. The algorithm
maintains original indices during elimination and records every pivot,
its complete restoring polynomial, and its sign. All arithmetic is
exact. A support cap chooses which substitutions to attempt; it never
truncates a polynomial or discards terms.

With support cap 64, 373 coordinates and residuals are removed from the
895-coordinate source, yielding 522 coordinates and 400 residuals. The
caps 16 and 4 yield the other two listed variants. No unused free
coordinate is discarded in these three releases, so the recorded
projection and restoration are actual bijections of natural zero sets.
Restoration is executed in reverse elimination order. Its composed
polynomial may have high degree; the final residuals still have degree
at most two, which is the degree property relevant to \eqref{eq:sos}.

The independent checker does not invoke the contraction routine or its
sparse substitution functions. It reconstructs each pivot from literal
coefficient lists, checks nonnegativity, expands each replacement, checks
degree after every step, and matches the final renaming and every final
residual. Exact squaring then certifies the delivered expanded polynomial.

\section{An unrestricted-integer companion without changing degree}
\label{sec:integers}

To remove ambiguity about the domain, we also deliver an equation over
all integers. For each coordinate $x_i$ of $\pn$, introduce three new
integer variables $a_i,b_i,c_i$ and define
\begin{equation}\label{eq:integer}
 \boxed{\displaystyle
 \pz(x,a,b,c)=\pn(x)+
 \sum_{i=0}^{521}(8x_i+1-a_i^2-b_i^2-c_i^2)^2=0.}
\end{equation}

At any integer zero, every guard vanishes, and hence $8x_i+1\ge0$.
As $x_i$ is integral, $x_i\ge0$. Conversely, for every $x_i\ge0$,
the positive integer $8x_i+1$ is congruent to one modulo eight and is
not of the forbidden form $4^a(8b+7)$. Legendre's three-square theorem
therefore supplies $a_i,b_i,c_i$ \cite{threesquares}. This proves the
existence equivalence between \eqref{eq:integer} over $\Z$ and
\eqref{eq:sos} over $\N$.

The integer polynomial is supplied in full, not just as a conversion
recipe. Its exact census is
\begin{center}
\begin{tabular}{lr}
\toprule
Unrestricted integer coordinates & 2,088\\
Quadratic residuals & 922\\
Expanded monomials & 26,791\\
Exact degree & 4\\
Maximum absolute coefficient & 83,901,150\\
\bottomrule
\end{tabular}
\end{center}
The new monomial $a_0^4$ has coefficient one, independently certifying
exact degree four. The conversion preserves existence; its fibers
generally contain several three-square representations, so it is not
claimed to be a bijection.

\section{Proof of the consistency theorem and its formalization scope}
\label{sec:scope}

We can now compose the argument without any assumption about whether
ZFC is consistent. From a hypothetical natural zero of $\pn$:
\begin{enumerate}
\item Restore every contracted coordinate and every removed affine gate.
      The sum of squares supplies the entire synchronized quadratic system.
\item Apply the common-frame and index theorems to recover every power.
      Recover every binomial by the common radix, then the factorials and
      modulus product by their exact graph identities.
\item Apply the shared CRT certificate to obtain every valid bounded
      local row. The local beta equations and endpoints form one path.
\item Apply the accumulator, carrier, and flow arguments to obtain a
      finite $C_{222}$ derivation of target code one.
\item Apply the unconditional calibrations \eqref{eq:calibration} to
      obtain a finite ZFC refutation.
\end{enumerate}
Conversely, a finite ZFC refutation gives a finite $C_{222}$ derivation
by the same calibration. Expand it to a tree, choose its label/schema
bound, form the carrier and state path, choose a beta encoding, and use
the CRT certificate. Evaluate the semantic arithmetic macros, choose
their common radix and shared Pell witnesses, and evaluate all gates.
Finally project along the recorded contractions. Every choice has been
proved to exist, with no bound on its size being assumed.

All these arguments use finite sequences, arithmetic induction, finite
CRT, elementary polynomial arithmetic, Pell classification and its
finite index criterion, and the stated calculus calibration. They can
be carried out in the internal natural numbers of ZFC. In particular,
there is no appeal to the arithmetic soundness of ZFC and no restriction
of its proof codes to externally standard numerals. This establishes
Theorem~\ref{thm:main} as an ordinary mathematical proof formalizable
in ZFC. A standard proof predicate for ZFC may be substituted for the
repository's sealed calculus using the usual effective proof translations.

The deliverable does \emph{not} contain a newly checked complete
Lean/Rocq derivation of this new polynomial theorem or an emitted
object-language ZFC proof term. The immediate inherited Lean statements
are included and pinned, but their entire transitive dependency closure
has not been rebuilt here. The exact Python audits establish finite
identities and artifact correspondence; they do not replace the universal
mathematical arguments above. They also do not prove that the equation
has no solution.

\section{Verification, reproduction, and integration}

The release separates universal proofs from finite checks. Its audits cover:
\begin{itemize}
\item exact unchanged constructor boundaries and the explicit factored
      Horn zero-set argument;
\item all 388 majorant inequalities, the literal 23-square construction,
      and positivity of every power exponent;
\item exact symbolic recovery of the synchronized power/index equations,
      the inclusion of all 22 squared ordinates in the common product,
      and the common binomial radix;
\item every contraction pivot, nonnegative restoring polynomial,
      intermediate degree, final coordinate map, and residual;
\item independent expansion and collection of the natural and integer
      polynomials, exact counts and coefficient heights;
\item deterministic regeneration and rejection of deliberately corrupted
      restoration coefficients, pivot signs, and coordinate maps.
\end{itemize}
Small Pell, CRT, binomial, and branch fixtures are supplementary regression
checks. No finite sample is presented as a proof of the unrestricted
consistency equivalence. Both ordinary Python and optimized Python
are used; authored validation checks raise explicit exceptions and do
not disappear when assertions are disabled.

The archive is self-contained for generation and the finite checks.
Python 3.10 or later and its standard library suffice. From its root:
\begin{verbatim}
python3 verify.py
python3 -O verify.py
python3 code/release.py --output /path/to/fresh-equations
sh build_report.sh
\end{verbatim}
The release generator writes a fresh equation tree. The verifier reads the
delivered tree, replays certificates, and regenerates in temporary staging;
it does not overwrite the mathematical data. The report is compiled from
the included LaTeX source with pdfLaTeX.

Integration can preserve this as a new contribution directory with the
report, equations, code, proofs, provenance, and receipts together. The
main comparison should record the 522-natural quartic and the
2,088-integer quartic separately. The 529- and 596-natural alternatives
belong alongside them as expression-size tradeoffs. The source snapshot
and the older calibration snapshot are distinct pins and are both recorded.

\section{What this iteration contributes}

The main change is a general method of organizing certificates across a
finite arithmetic circuit. Large bounds are chosen for a whole family;
Pell and CRT witnesses are chosen compatibly for that family; and
semantic values are separated from auxiliary encoding choices. This
eliminates repeated local certificates without weakening their arithmetic
content. The small-coefficient, fully expanded direct equation makes the
gain concrete: 522 rather than 3,614 natural coordinates, with degree
four retained and every coefficient supplied.

The next substantial reduction would have to simplify more of the direct
calculus or replace the remaining Pell-index and arithmetic-graph
structures. The present contraction order is not proved optimal. The
current result already supplies a reusable simultaneous-power theorem,
a common-radix binomial theorem, and a complete finite polynomial whose
logical meaning is precisely stated. It is a measured improvement over
the inspected direct construction, without claiming to resolve the
unrestricted search for a globally simplest consistency equation.

\begin{thebibliography}{99}
\bibitem{repository} V. Reshetnikov and ProveIt contributors.
\emph{DiophantineConsistency research}, snapshot
\file{a88aefd7ddab4771534bb049c67350e397c91b9a}.
\url{https://github.com/VladimirReshetnikov/ProveIt/tree/a88aefd7ddab4771534bb049c67350e397c91b9a/Computability/DiophantineConsistency}.

\bibitem{latest} ProveIt contributors.
\emph{Later arrivals and current size reductions},
\file{PROPOSAL_REVIEW_LATEST.md}, same snapshot. An exact local copy is
included under \file{provenance/repository}.

\bibitem{shared} ProveIt incoming contribution.
\emph{Direct certificates and coordinated Diophantine reductions},
\file{ConZFC_Shared_Certificates_2026-10-24.zip}, as present in the
comparison snapshot. Immediate arithmetic source, source SHA-256 values,
and original notices are retained under \file{provenance/upstream}.

\bibitem{calibration} ProveIt contributors.
\file{Code254.lean}, \file{KCode222.lean}, and \file{KEndpoint.lean},
calibration snapshot \file{08200eb10fde2db408c69dffd840d6920a569b1f}.
\url{https://github.com/VladimirReshetnikov/ProveIt/blob/08200eb10fde2db408c69dffd840d6920a569b1f/Computability/BusyBeaver/Lean/BusyBeaver/ZFC/KEndpoint.lean}.

\bibitem{choice} ProveIt contributors.
\file{ConstructibleChoice.lean}, same calibration snapshot.
\url{https://github.com/VladimirReshetnikov/ProveIt/blob/08200eb10fde2db408c69dffd840d6920a569b1f/Computability/BusyBeaver/Lean/BusyBeaver/ZFC/ConstructibleChoice.lean}.

\bibitem{pell} Mathlib contributors.
\emph{Pell's equation and Matiyasevic's theorem},
\file{Mathlib.NumberTheory.PellMatiyasevic}, especially
\file{eq_pell}, \file{xn_ge_a_pow}, \file{yn_modEq_a_sub_one},
and \file{matiyasevic}.
\url{https://leanprover-community.github.io/mathlib4_docs/Mathlib/NumberTheory/PellMatiyasevic.html}.

\bibitem{carneiro} M. Carneiro.
\emph{A Lean formalization of Matiyasevi\v{c}'s theorem}, 2018.
\url{https://arxiv.org/abs/1802.01795}.

\bibitem{pellinterface} ProveIt contributors.
\file{PellRepresentation.lean}, especially
\file{criterion_iff_pell}, calibration snapshot above. Exact source is
included in \file{provenance/upstream/calibration}.

\bibitem{jones} J. P. Jones.
Universal Diophantine equation.
\emph{Journal of Symbolic Logic} \textbf{47} (1982), 549--571.
\url{https://doi.org/10.2307/2273588}.

\bibitem{sun} Z.-W. Sun.
\emph{Ten unknowns for Hilbert's tenth problem over the integers},
2026 preprint.
\url{https://arxiv.org/abs/2609.01594}.

\bibitem{threesquares} P. Pollack and E. Trevi\~no.
Dirichlet's proof of the three-square theorem: an algorithmic perspective.
\emph{Mathematics of Computation} \textbf{88} (2019).
\url{https://www.ams.org/mcom/2019-88-316/S0025-5718-2018-03349-X/}.
We use only the classical unconditional three-square theorem, specialized
to $8n+1$, not any conditional running-time bound for finding a representation.
\end{thebibliography}
\end{document}
